
\section{Relaciones con la completación}
\label{IV.3}

Sean~$A$ un anillo noetheriano, $I$ un ideal de~$A$,
$\widehat{A}$ la completación separada de~$A$ para la topología
$I$-preádica y, para todo $A$-módulo~$M$, sea~$\widehat{M}$ su
completación para la topología $I$-preádica. Este es un
$\widehat{A}$-módulo, de donde se obtiene un homomorfismo canónico
$$
M\otimes_A\widehat{A}\to\widehat{M}.
$$
Cuando~$M$ recorre los módulos \emph{de tipo finito}, el funtor
$M\mto\widehat{M}$ es exacto, como resulta fácilmente del
\emph{teorema de Krull: si~$N\subset M$, la topología de~$N$ es la
inducida por la topología de~$M$}. Como el funtor
$M\mto M\otimes_A\widehat{A}$ es exacto por la derecha, se concluye
fácilmente (resolviendo~$M$ mediante~$L\to L'\to M$, con~$L$ y~$L'$
libres de tipo finito) que el homomorfismo funtorial anterior es un
\emph{isomorfismo} (pues el funtor de completación también es exacto por
la derecha) y, por consiguiente, que~$M\mto M\otimes_A\widehat{A}$ es
también un funtor \emph{exacto}. Por tanto:

\begin{proposition}
\label{IV.3.1}
Sean~$A$ un anillo noetheriano e~$I$ un ideal de~$A$. Entonces la
completación separada~$\widehat{A}$ de~$A$ para la topología
$I$-preádica es \emph{plana} sobre~$A$.
\end{proposition}

\begin{corollaire}
\label{IV.3.2}
Para que~$\widehat{A}$ sea fielmente plano sobre~$A$, es necesario y
suficiente que~$I$ esté contenido en el radical de Jacobson de~$A$.
\end{corollaire}

En efecto, basta aplicar el criterio~\Ref{IV.2.1}~(i~ter).

Estos resultados resumen todo cuanto puede decirse, desde el punto de vista
del álgebra lineal, sobre las relaciones entre~$A$ y~$\widehat{A}$. El
corolario~\Ref{IV.3.2} se utiliza sobre todo cuando~$A$ es un anillo local
noetheriano e~$I$ está contenido en el ideal maximal~$\goth{m}$ (y, casi
siempre, es igual a él).

